% ------------------------------------------------------------------
%  Capitolo 2 — Come si studia all'università, oggi
%  Parte I
%  Ricerca di riferimento: ricerca 02 §2–6, §9; 01 §1.2
%  Voci bibliografiche principali: karpicke2009, hartwig2012, kornell2007, almalaurea2026, oecd2026pisa, invalsi2026
%  Epigrafe: ✔ verificata sul testo (vedi ricerca 03 e registro, ricerca 00)
%  Scheda del capitolo: ricerca/06_indice-ragionato.md, capitolo 2
% ------------------------------------------------------------------
\chapter{Come si studia all'università, oggi}
\label{cap:oggi}

\epigrafe[sopporti all'inizio la noia di rileggere più volte ciò che ha scritto e letto]{devoret initio taedium illud et scripta et lecta saepius revolvendi}{Quintiliano, Institutio oratoria XI, 2, 41}

\iniziale{Q}{uesto capitolo} \segnaposto{apertura: da scrivere — vedi la scheda nell'indice ragionato}.

\section{Rileggere, sottolineare, ripetere}

\segnaposto{da scrivere}

\section{Che cosa ci si porta dalla scuola}

\segnaposto{da scrivere}

\section{I numeri dell'università italiana}

\segnaposto{da scrivere}

\section{Nessuno ci ha insegnato a studiare}

\segnaposto{da scrivere}

\begin{daricordare}
\segnaposto{il principio del capitolo in due righe}
\end{daricordare}

\begin{domande}
  \item \segnaposto{domanda di recupero su questo capitolo}
  \item \segnaposto{domanda di applicazione a una materia che studi}
  \item \segnaposto{domanda di ripresa su un capitolo precedente}
\end{domande}

\begin{sintesi}
  \item \segnaposto{punto di sintesi}
\end{sintesi}

\finecapitolo
